\documentclass[a4paper]{article}     
\usepackage{amsfonts}
\usepackage{amsmath}
\usepackage{amssymb}
\usepackage{graphicx}
\setlength{\parindent}{0pt}

\newcommand{\fx}{\dfrac{x^2-5x+6}{x^3-2x^2-x+2}}
\newcommand{\llim}[1]{ \lim\limits_{x \rightarrow #1} }
\newcommand{\fxf}{\dfrac{(x-3)(x-2)}{(x-2)(x^2-1)}}

\begin{document}

\noindent \large Auteur: Milo Rampelbergh \\
\noindent \large Studierichting: Informatica\\
\noindent \large Oefeningenreeks 2D nr. 6.2\\

\medskip

\normalsize

%Begin oefening

$ f(x) = \fx $\\

\begin{enumerate}
\item $ \llim{2} \fx = \dfrac{0}{0} $\\

\hspace{1.5cm} $ x^2-5x+6 = 0 \rightarrow x = 3 $ of $ x = 2 $ \\

\hspace{1.5cm} $ \begin{array}{c|ccc|c}
& 1 & -2 & -1 & 2 \\
2 & & 2 & 0 & -2 \\
\hline
& 1 & 0 & -1 & 0 \\
\end{array} $ \\

$ = \llim{2} \dfrac{(x-3)(x-2)}{(x-2)(x^2-1)} $ \\

$ = \llim{2} \dfrac{(x-3)}{(x^2-1)} = \dfrac{-1}{3} $ \\

\item $ \llim{3} \fx = \dfrac{-15+6}{27-18-3+2} = \dfrac{0}{2} = 0 $ \\

\item $ \llim{-1} \fx = \dfrac{1+5+6}{-1-2+1+2} = \dfrac{12}{0} = \infty $ \\

$ \begin{array}{ll}
\Rightarrow & \llim{-1+} \fx = \llim{-1+} \fxf = \dfrac{12}{0^+} = +\infty \\
& \llim{-1-} \fx = \llim{-1-} \fxf = \dfrac{12}{0^-} = -\infty \\
\end{array} $  \\

\item $ \llim{\infty} \fx = \llim{\infty} \dfrac{x^2}{x^3} = \llim{\infty} \dfrac{1}{x} = 0 $ \\

\item $ \llim{1} \fx = \dfrac{1-5+6}{1-2-1+2} = \dfrac{2}{0} = \infty $ \\

$ \begin{array}{ll}
\Rightarrow & \llim{1+} \fx = \llim{1+} \fxf = \dfrac{2}{0^+} = -\infty \\
& \llim{1-} \fx = \llim{1-} \fxf = \dfrac{2}{0^-} = +\infty \\
\end{array} $  \\

\end{enumerate}

%Einde oefening

\begin{thebibliography}{999}
%\bibitem{a} Referentie
\end{thebibliography}
\end{document} 